\documentclass[11pt,a4paper]{moderncv}

% Estilo y color
\moderncvstyle{classic}
\moderncvcolor{blue}

% Ajuste de márgenes para asegurar que quepa en una página
\usepackage[scale=0.88]{geometry}

% Datos personales
\name{Andrés}{López}
\title{Data Scientist \& Físico}
\address{Ciudad de México, México}{}{}
\phone[mobile]{+52~1~55~6543~4108}
\email{anderiotti@ciencias.unam.mx}
\social[linkedin]{Andres Lopez}

\begin{document}
\makecvtitle

\section{Perfil}
Estudiante avanzado de Física y Ciencia de Datos con orientación en Quantum Machine Learning, combinando un riguroso perfil analítico con la disciplina y resiliencia del alto rendimiento como miembro del equipo representativo de powerlifting de la UNAM. Apasionado por transformar conocimiento científico en modelos predictivos de alto impacto mediante la integración de modelado matemático, análisis estadístico, implementación computacional de algoritmos y optimización.

\section{Educación}
\cventry{2021--Presente}{Licenciatura en Física}{Universidad Nacional Autónoma de México}{CDMX}{}{
Enfoque en modelado de sistemas complejos, mecánica cuántica y análisis numérico. Experiencia en machine learning en laboratorios de física experimental. \textbf{Promedio: 9.5/10.}}

\cventry{2021--Presente}{Licenciatura en Ciencia de Datos}{Universidad Nacional Autónoma de México}{CDMX}{}{
Especialización en optimización, teoría de grafos y estadística multivariada. Desarrollo de modelos de regresión, clasificación y clustering. \textbf{Promedio: 9.6/10.}}

\section{Experiencia Relevante}
\cventry{2024--2025}{Analista de Datos}{Grupo de Investigación Estudiantil, UNAM}{CDMX}{}{
Desarrollo de modelos predictivos en \textbf{Python} y \textbf{R} para series temporales de experimentos físicos. Limpieza de datos con \emph{SQL} y planificación de sistemas OLTP.}

\cventry{2025}{(IBM) Qiskit Global Summer School}{Quantum Badge}{}{}{
Desarrollo de modelos de Quantum Machine Learning (QML) para clasificación de datos y mitigación de errores, incluyendo programación en hardware cuántico real de IBM.}

\cventry{2024}{Proyecto de Análisis Magnetohidrodinámico}{Lab. de Hidrodinámica y Turbulencia, UNAM}{CDMX}{}{
Lideré el modelado de la transición de un fluido cargado predecible a un régimen caótico, cuantificando el impacto de la topología del campo electromagnético en la dinámica del sistema.}

\section{Competencias Técnicas}
\cvitem{Lenguajes}{Python, Julia, R, C, C++, SQL}
\cvitem{Machine Learning}{Scikit-learn, TensorFlow (Keras), PyTorch, clasificación, regresión, series temporales}
\cvitem{Quantum Computing}{Qiskit, PennyLane, fundamentos de computación cuántica, modelos híbridos QML}
\cvitem{Data Analysis}{Pandas, NumPy, Matplotlib, Seaborn, análisis exploratorio, optimización}
\cvitem{Business Intelligence}{Power BI, Tableau, Excel (PowerQuery, tablas dinámicas), Power Designer}
\cvitem{Bases de Datos}{Diseño y consultas en SQL, modelado relacional, SQLAlchemy}
\cvitem{Nube y Dev Tools}{Git, Linux, Docker básico, AWS, VirtualBox}

\section{Idiomas}
\cvitemwithcomment{Inglés}{Avanzado}{Lectura y conversación fluida}

\end{document}